% Options for packages loaded elsewhere
% Options for packages loaded elsewhere
\PassOptionsToPackage{unicode}{hyperref}
\PassOptionsToPackage{hyphens}{url}
\PassOptionsToPackage{dvipsnames,svgnames,x11names}{xcolor}
%
\documentclass[
  a4paper,
  DIV=11,
  numbers=noendperiod]{scrartcl}
\usepackage{xcolor}
\usepackage[margin=2cm]{geometry}
\usepackage{amsmath,amssymb}
\setcounter{secnumdepth}{-\maxdimen} % remove section numbering
\usepackage{iftex}
\ifPDFTeX
  \usepackage[T1]{fontenc}
  \usepackage[utf8]{inputenc}
  \usepackage{textcomp} % provide euro and other symbols
\else % if luatex or xetex
  \usepackage{unicode-math} % this also loads fontspec
  \defaultfontfeatures{Scale=MatchLowercase}
  \defaultfontfeatures[\rmfamily]{Ligatures=TeX,Scale=1}
\fi
\usepackage{lmodern}
\ifPDFTeX\else
  % xetex/luatex font selection
\fi
% Use upquote if available, for straight quotes in verbatim environments
\IfFileExists{upquote.sty}{\usepackage{upquote}}{}
\IfFileExists{microtype.sty}{% use microtype if available
  \usepackage[]{microtype}
  \UseMicrotypeSet[protrusion]{basicmath} % disable protrusion for tt fonts
}{}
\makeatletter
\@ifundefined{KOMAClassName}{% if non-KOMA class
  \IfFileExists{parskip.sty}{%
    \usepackage{parskip}
  }{% else
    \setlength{\parindent}{0pt}
    \setlength{\parskip}{6pt plus 2pt minus 1pt}}
}{% if KOMA class
  \KOMAoptions{parskip=half}}
\makeatother
% Make \paragraph and \subparagraph free-standing
\makeatletter
\ifx\paragraph\undefined\else
  \let\oldparagraph\paragraph
  \renewcommand{\paragraph}{
    \@ifstar
      \xxxParagraphStar
      \xxxParagraphNoStar
  }
  \newcommand{\xxxParagraphStar}[1]{\oldparagraph*{#1}\mbox{}}
  \newcommand{\xxxParagraphNoStar}[1]{\oldparagraph{#1}\mbox{}}
\fi
\ifx\subparagraph\undefined\else
  \let\oldsubparagraph\subparagraph
  \renewcommand{\subparagraph}{
    \@ifstar
      \xxxSubParagraphStar
      \xxxSubParagraphNoStar
  }
  \newcommand{\xxxSubParagraphStar}[1]{\oldsubparagraph*{#1}\mbox{}}
  \newcommand{\xxxSubParagraphNoStar}[1]{\oldsubparagraph{#1}\mbox{}}
\fi
\makeatother


\usepackage{longtable,booktabs,array}
\usepackage{calc} % for calculating minipage widths
% Correct order of tables after \paragraph or \subparagraph
\usepackage{etoolbox}
\makeatletter
\patchcmd\longtable{\par}{\if@noskipsec\mbox{}\fi\par}{}{}
\makeatother
% Allow footnotes in longtable head/foot
\IfFileExists{footnotehyper.sty}{\usepackage{footnotehyper}}{\usepackage{footnote}}
\makesavenoteenv{longtable}
\usepackage{graphicx}
\makeatletter
\newsavebox\pandoc@box
\newcommand*\pandocbounded[1]{% scales image to fit in text height/width
  \sbox\pandoc@box{#1}%
  \Gscale@div\@tempa{\textheight}{\dimexpr\ht\pandoc@box+\dp\pandoc@box\relax}%
  \Gscale@div\@tempb{\linewidth}{\wd\pandoc@box}%
  \ifdim\@tempb\p@<\@tempa\p@\let\@tempa\@tempb\fi% select the smaller of both
  \ifdim\@tempa\p@<\p@\scalebox{\@tempa}{\usebox\pandoc@box}%
  \else\usebox{\pandoc@box}%
  \fi%
}
% Set default figure placement to htbp
\def\fps@figure{htbp}
\makeatother





\setlength{\emergencystretch}{3em} % prevent overfull lines

\providecommand{\tightlist}{%
  \setlength{\itemsep}{0pt}\setlength{\parskip}{0pt}}



 


\KOMAoption{captions}{tableheading}
\makeatletter
\@ifpackageloaded{caption}{}{\usepackage{caption}}
\AtBeginDocument{%
\ifdefined\contentsname
  \renewcommand*\contentsname{Table of contents}
\else
  \newcommand\contentsname{Table of contents}
\fi
\ifdefined\listfigurename
  \renewcommand*\listfigurename{List of Figures}
\else
  \newcommand\listfigurename{List of Figures}
\fi
\ifdefined\listtablename
  \renewcommand*\listtablename{List of Tables}
\else
  \newcommand\listtablename{List of Tables}
\fi
\ifdefined\figurename
  \renewcommand*\figurename{Figure}
\else
  \newcommand\figurename{Figure}
\fi
\ifdefined\tablename
  \renewcommand*\tablename{Table}
\else
  \newcommand\tablename{Table}
\fi
}
\@ifpackageloaded{float}{}{\usepackage{float}}
\floatstyle{ruled}
\@ifundefined{c@chapter}{\newfloat{codelisting}{h}{lop}}{\newfloat{codelisting}{h}{lop}[chapter]}
\floatname{codelisting}{Listing}
\newcommand*\listoflistings{\listof{codelisting}{List of Listings}}
\makeatother
\makeatletter
\makeatother
\makeatletter
\@ifpackageloaded{caption}{}{\usepackage{caption}}
\@ifpackageloaded{subcaption}{}{\usepackage{subcaption}}
\makeatother
\usepackage{bookmark}
\IfFileExists{xurl.sty}{\usepackage{xurl}}{} % add URL line breaks if available
\urlstyle{same}
\hypersetup{
  pdftitle={Tree based methods},
  colorlinks=true,
  linkcolor={blue},
  filecolor={Maroon},
  citecolor={Blue},
  urlcolor={Blue},
  pdfcreator={LaTeX via pandoc}}


\title{Tree based methods}
\author{}
\date{}
\begin{document}
\maketitle


\subsection{Part 1 - Getting to know the
data}\label{part-1---getting-to-know-the-data}

Run the python code in Part 1 and answer the following questions

1.1 How many predictors and how many data points does the dataset have?

\begin{center}\rule{0.5\linewidth}{0.5pt}\end{center}

1.2 The response variable is \texttt{Class}, how many categories does it
have?Is the dataset balanced?

\begin{center}\rule{0.5\linewidth}{0.5pt}\end{center}

1.3 Why does it matter if the dataset is balenced or not? What could
happen with a very unbalanced dataset?

\begin{center}\rule{0.5\linewidth}{0.5pt}\end{center}

1.4 Looking at the generated plots, which variables seem most useful?

\begin{center}\rule{0.5\linewidth}{0.5pt}\end{center}

\subsection{Part 2: A first, simple
tree}\label{part-2-a-first-simple-tree}

2.1 Which variable is used to split the data? What is the threshold
value used in the split and what does this threshold mean in practice?

\begin{center}\rule{0.5\linewidth}{0.5pt}\end{center}

2.2 How would you explain this model to someone with no statistics
background?

\begin{center}\rule{0.5\linewidth}{0.5pt}\end{center}

2.3 A bean has Area = 50,000. What is the predicted bean variety?
Another bean has Area = 80,000. What is the predicted bean variety? What
happens if two beans have the same value of Area but very different
values of Compactness?

\begin{center}\rule{0.5\linewidth}{0.5pt}\end{center}

\subsection{2.4 How, do you think, the tree produces prediction for each
group?}\label{how-do-you-think-the-tree-produces-prediction-for-each-group}

2.5 Do you think that the model is overfitting or underfitting?

\begin{center}\rule{0.5\linewidth}{0.5pt}\end{center}

\subsection{Part 3: A deeper tree}\label{part-3-a-deeper-tree}

3.1 Predict the class for the following beans based on the tree with
depth = 6 created in Part 3 of the code:

\subsubsection{Bean A}\label{bean-a}

\begin{itemize}
\tightlist
\item
  \textbf{Area:} 40,000
\item
  \textbf{Perimeter:} 800
\item
  \textbf{MajorAxisLength:} 300
\item
  \textbf{MinorAxisLength:} 180
\item
  \textbf{Compactness:} 0.75
\end{itemize}

\textbf{Predicted variety:} \_\_\_\_\_\_\_\_\_\_\_\_\_\_

\subsubsection{Bean B}\label{bean-b}

\begin{itemize}
\tightlist
\item
  \textbf{Area:} 90,000
\item
  \textbf{Perimeter:} 1200
\item
  \textbf{MajorAxisLength:} 450
\item
  \textbf{MinorAxisLength:} 250
\item
  \textbf{Compactness:} 0.85
\end{itemize}

\textbf{Predicted variety:} \_\_\_\_\_\_\_\_\_\_\_\_\_\_\_\_\_\_\_\_

\textbf{Predicted species:} \_\_\_\_\_\_\_\_\_\_\_\_\_\_\_\_\_\_\_\_

3.2 Which variable appears most often in the splitting rules? What does
this suggest about its usefulness for classifying bean species?

\begin{center}\rule{0.5\linewidth}{0.5pt}\end{center}

3.3 Do you think this tree can be generalized on new bean data? Explain
your reasoning and discuss whether the tree might be overfitting.

\begin{center}\rule{0.5\linewidth}{0.5pt}\end{center}

\subsection{Part 4 - Tuning the tree
depth}\label{part-4---tuning-the-tree-depth}

In part 4 the code performs a cross-validation experiment to tune the
tree depth parameter.

4.1 How many folds are used in the cross-validation procedure? How can
you tell from the code?

\begin{center}\rule{0.5\linewidth}{0.5pt}\end{center}

The table produced reports the mean accuracy for each value of the tree
depth. The accuracy is defined as: \[
\text{Accuracy} = \frac{\text{Number of correct predictions}}{\text{Total numnef of predictions}}
\]

4.2 Why does the output report the ``Mean CV Accuracy'' rather than a
single accuracy value? What does the reported standard deviation tell
us?

\begin{center}\rule{0.5\linewidth}{0.5pt}\end{center}

4.3 Does performance improve indefinitely as depth increases? Why might
very deep trees perform worse than shallow tree? Which depth value would
you choose?

\begin{center}\rule{0.5\linewidth}{0.5pt}\end{center}

\subsection{Part 5 - Fitting a Randon
Forest}\label{part-5---fitting-a-randon-forest}

Run part of the Python code. Here we fit a random forest to the beans
data.

5.1 How many trees are included in this random forest? How can you tell
from the code?

\begin{center}\rule{0.5\linewidth}{0.5pt}\end{center}

5.2 Before evaluating the model, would you expect the random forest to
perform better or worse than a single decision tree? Explain your
reasoning.

\begin{center}\rule{0.5\linewidth}{0.5pt}\end{center}

5.3 What is one advantage and one disadvantage of a random forest
compared with a single decision tree?

\begin{center}\rule{0.5\linewidth}{0.5pt}\end{center}

\subsection{Part 6: Cross validation for random
forests}\label{part-6-cross-validation-for-random-forests}

Run the code in Part 6. This is a 5-fold cross validation to tune two of
the main parameters in a random forest: the tree depth and the number of
trees in the forest.

Note: the code takes some seconds to be finished.

6.1 Based on these results, what values of \texttt{n\_estimators} and
\texttt{max\_depth} would you choose for the final model?

\begin{center}\rule{0.5\linewidth}{0.5pt}\end{center}

\subsection{Part 7: Testing on the test
set}\label{part-7-testing-on-the-test-set}

7.1 If your goal were prediction, which model would you choose: the best
decision tree or the best random forest? Explain your answer using the
test-set results.

\begin{center}\rule{0.5\linewidth}{0.5pt}\end{center}

7.2 What is the accuracy of the selected model on the test set? How does
this compare with the cross-validation accuracy obtained during model
selection?

\begin{center}\rule{0.5\linewidth}{0.5pt}\end{center}

7.3 If your goal were interpretability, which model would you choose?
Why?

\begin{center}\rule{0.5\linewidth}{0.5pt}\end{center}

7.4 Examine the confusion matrix.

\begin{itemize}
\tightlist
\item
  Which bean variety is most often classified correctly?
\item
  Which bean variety is most often misclassified?
\item
  Which pair of bean varieties appears most difficult to distinguish?
\end{itemize}

\begin{center}\rule{0.5\linewidth}{0.5pt}\end{center}




\end{document}
